% =====================================================================
% §2 CONCEPTUAL SECTION — new material
% Pattern: economic problem -> theory in prose, no math -> assumptions
%          enumerated -> closes on the empirical question
% See skills/realestate-finance-writing/references/structure.md
%
% WIRING DEPENDENCY: this file contains fig:equilibrium_outcomes, moved
% here from §3.1. Deleting the float from §3.1 is required in the same
% commit that inputs this file, or the label will collide.
% =====================================================================

\section{Unraveling with Possibly Uninformed Sellers}
\label{sec:unraveling}

\subsection{The Unraveling Benchmark}

Asymmetric information about product quality can prevent trade from happening at the right
prices. If buyers cannot verify quality before purchasing, they bid no more than the value of
what they expect to receive on average. Sellers of genuinely good products are then underpaid,
and may withdraw from the market altogether, leaving buyers facing a pool of worse quality than
they had hoped for. The outcome is bad for both sides: buyers pay for uncertainty they would
rather not bear, and the best sellers cannot capture the value of what they own.

Disclosure appears to resolve this problem without any regulatory intervention. Suppose quality
is verifiable, so that a seller who makes a claim about it can be held to that claim. A seller
of the highest-quality product will disclose, because staying silent means being valued at the
average of the pool rather than at what the product is actually worth. Buyers, anticipating
this, understand that any seller who remains silent is not of the highest quality, and revise
downward what they will pay for silence. That revision gives the second-highest-quality seller
the same incentive, and disclosing becomes worthwhile for them too. The argument repeats. Each
round of disclosure worsens the expected quality of the remaining silent pool, which strengthens
the incentive for the best of those who remain to separate themselves. The process continues
until all sellers but possibly the very worst have disclosed. This is the unraveling result of
\citet{grossman_informational_1981} and \citet{milgrom_good_1981}: even sellers with bad news
disclose, because they anticipate that silence will be read as worse news still. If it holds,
market forces alone deliver near-complete disclosure and a mandate is redundant.

The argument is powerful but it is not free. It rests on four assumptions. First, (a)~buyers know
that sellers possess private information about quality; if a seller may simply not know, silence
no longer implies unwillingness to reveal. Second, (b)~disclosure is costless, or at least cheap
relative to the price consequences of silence; a large literature examines how disclosure costs
truncate the unraveling process and leave a residual pool of non-disclosers
\citep{verrecchia_1983, jovanovic_1982}. Third, (c)~disclosure is credible, so that a claim about
quality cannot be fabricated; without verifiability, disclosure conveys nothing and the
unraveling logic never starts. Fourth, (d)~buyers are homogeneous in how much they value the
disclosed attribute, so that the price penalty for silence is felt by every seller with equal
force.

The first three assumptions have received sustained attention. The fourth has not, even though
it governs the strength of the entire mechanism: what drives a seller to disclose is the revenue
lost when buyers price silence pessimistically, and that loss is proportional to the share of
buyers who care about the attribute in the first place. Whether unraveling describes any
particular market is therefore an empirical question, and answering it requires knowing which of
these assumptions fails there.

\subsection{Two Departures and Why They Compound}

Two of the four assumptions are implausible in markets where quality is certified by a
third party and buyers differ in their exposure to the certified attribute.

The first departure concerns assumption~(a). When some sellers genuinely lack private
information, silence pools those who \emph{will not} disclose with those who \emph{cannot}.
\citet{dye_disclosure_1985} showed that this possibility breaks the buyer inference on which
unraveling depends: a rational buyer confronted with silence must now form a posterior over two
distinct populations, and how pessimistic that posterior should be depends on their relative
masses rather than on the act of silence itself. \citet{jung_kwon_1988} characterised the
resulting equilibria and showed that partial pooling, in which some informed sellers of poor
quality withhold while the uninformed are silent by necessity, is generically stable. Disclosure
does not unravel to the bottom; it stops somewhere in the interior.

The second departure concerns assumption~(d). If only a fraction of buyers price the disclosed
attribute, a seller who withholds it faces a diluted penalty, because part of the demand side is
indifferent to what was withheld. \citet{fishman_mandatory_2003} show formally that mandated
disclosure can be ineffective, and even welfare-reducing, when buyers differ in their capacity to
use disclosed information. The disclosure incentive is thus scaled by the composition of demand,
and where that composition is unfavourable the incentive can be too weak to sustain separation
even among sellers who are perfectly well informed.

The two departures do not merely coexist; they compound. A larger mass of sellers who cannot
disclose dilutes the adverse inference buyers draw from silence, which raises the expected quality
of the silent pool and makes withholding more attractive to informed sellers of poor quality.
Those sellers then join the pool, and while their entry worsens its true average quality, buyers
cannot separate them from the involuntarily silent. Meanwhile, a demand side in which many buyers
do not price the attribute weakens the penalty that would otherwise discipline them. Put together,
the consequence is stronger than persistence of non-disclosure: silence can command a
\emph{premium} over the worst disclosed quality, because the silent pool contains a substantial
mass of sellers whose silence carries no adverse information at all. The sellers who cannot
disclose do not intend to shelter anyone. Their presence is involuntary. But it is precisely what
makes silence ambiguous, and ambiguity is what makes concealment privately rational.

Because the two departures operate on different primitives, they map out a space of possible
market outcomes rather than a single prediction. Figure~\ref{fig:equilibrium_outcomes} organises
that space along the two dimensions: the share of sellers who hold verifiable information, and
the share of buyers who price it. Full unraveling requires both to be high. Where either is
interior, silence survives in equilibrium, and where the informed share is low while buyer
sensitivity is not, silence can be rewarded rather than punished.

\begin{figure}[H]
    \centering
    \includegraphics[width=\textwidth]{fig/visuals/equilibirum_outcomes.png}
    \caption{Equilibrium outcomes under seller information and buyer sensitivity heterogeneity}
    \label{fig:equilibrium_outcomes}
\end{figure}

Two features of this taxonomy matter for what follows. It implies that the informational content
of silence is a market-level object rather than a universal one, so that two markets facing the
same disclosure mandate can sit in different cells and sustain opposite price patterns. And it
implies that the two primitives are separately movable by policy: the share of sellers holding
verifiable information can be raised by investment in certification infrastructure, while the
share of buyers who price the attribute responds to the legal weight the certificate carries and
to the cost of the attribute itself. Which cell a given market occupies, however, cannot be
settled by theory. It requires measuring the composition of the silent pool and the sensitivity
of buyers separately, which is what the data described in the next section allow us to do.
